\section{Proofs for the estimated-matrix comparison}
\label{app:multivariate-waiting}
All limits and moment estimates in this appendix are under $P$ unless marked
$Q$. Dimension, shrinkage, positive ridge, leverage and training size are fixed.

\subsection{Integrated estimation and moments}
Let $S_m$ be partial sums of iid standard $p$-dimensional Gaussian vectors.
For any fixed symmetric $B\succeq0$ and $i\le j$, Gaussian fourth moments give
\[
 \E(S_m'BS_m)=m\operatorname{tr}B,
 \qquad
 \operatorname{Cov}(S_i'BS_i,S_j'BS_j)
       =2i^2\operatorname{tr}(B^2).
\]
Therefore $\sum_{m\le n}S_m'BS_m/m^2$ has mean
$\operatorname{tr}B\log n+O(1)$ and variance $O(\log n)$.
The positive-summand subsequence argument in
Appendix~\ref{app:vanishing-normalization} yields
\[
 \frac1{\log n}\sum_{m\le n}\frac{S_m'BS_m}{m^2}
       \longrightarrow\operatorname{tr}B\quad\text{almost surely}.
\]
Apply this with
$B=\Sigma^{1/2}H^{-1}\Sigma H^{-1}\Sigma^{1/2}$, whose trace is $\chi$,
and with $B=\Sigma$. The ordinary strong laws give
$\widehat H_t\to H$ and $\widehat\mu_t\to0$. Thus the unconstrained
optimizer tends to zero and is eventually feasible for the $L_1$ constraint.
For those dates, the difference between $v_t$ and
$\widehat\mu_t'H^{-1}\Sigma H^{-1}\widehat\mu_t/\gamma^2$ is bounded by
$\epsilon_t\norm{\widehat\mu_t}^2$, with $\epsilon_t\to0$ almost surely.
The second integrated law and weighted Cesaro imply that their accumulated
difference divided by $\log n$ vanishes. Removing the fixed training prefix
proves $\sum_{t\le n}v_t/\log n\to\chi/\gamma^2$.

The first-order condition along the feasible direction $-w_t$ gives
\[
 \gamma r\norm{w_t}_2^2
 \le\gamma w_t'\widehat H_tw_t
 \le\widehat\mu_t'w_t,
 \qquad \norm{w_t}_2\le\frac{\norm{\widehat\mu_t}_2}{\gamma r}.
\]
Since $\widehat\mu_t\sim N_p(0,\Sigma/m_t)$, for each fixed positive $a$,
$\E\norm{w_t}_2^a=O(m_t^{-a/2})$. In particular
\[
 \E v_t=O(m_t^{-1}),\qquad
 \sum_t\E v_t^{3/2}<\infty,\qquad
 \sum_t\E v_t^2<\infty,
\]
and $v_t\le L^2\norm\Sigma_{\op}$ deterministically. More explicitly,
\[
 \E v_t^2\le
 \frac{\norm\Sigma_{\op}^2
       \{(\operatorname{tr}\Sigma)^2+2\operatorname{tr}(\Sigma^2)\}}
      {\gamma^4r^4m_t^2}.
\]
These inequalities require no inverse-Wishart moment or growing-dimension
approximation. They also apply when $\Sigma$ is singular.

\subsection{A rank-one conditional null}
Put $a_t=\gamma\sqrt{v_t}/2$. When $v_t>0$, define the conditional Gaussian
kernel under $Q$ by mean and covariance
\[
 \mu_{Q,t}=\frac{\gamma\Sigma w_t}{2(1+a_t^2)},\qquad
 \Sigma_{Q,t}=\Sigma-
 \{1-(1+a_t^2)^{-2}\}\frac{\Sigma w_tw_t'\Sigma}{v_t}.
\]
Use $N_p(0,\Sigma)$ when $v_t=0$, and the original training law. In whitened
coordinates this changes only the unit direction
$e_t=\Sigma^{1/2}w_t/\sqrt{v_t}$: its mean is $a_t/(1+a_t^2)$ and variance
$(1+a_t^2)^{-2}$; orthogonal directions retain their original centered
Gaussian law. For singular $\Sigma$, work on its range. These kernels have
the same support and strictly positive densities there at every finite date.

Their raw second-moment matrix is
\[
 M_{Q,t}=\Sigma-\frac{a_t^2}{1+a_t^2}
                    \frac{\Sigma w_tw_t'\Sigma}{v_t}\preceq\Sigma.
\]
Also $w_t'\mu_{Q,t}=\gamma v_t/\{2(1+a_t^2)\}$ and
$w_t'M_{Q,t}w_t=v_t/(1+a_t^2)$, proving $\E_Q(Y_t\mid\F_{t-1})=0$.
For each direction $b$, a Gaussian variable with mean $\mu$ and raw second
moment $M$ has fourth moment $3M^2-2\mu^4$. Consequently the displayed
directional fourth-moment contract follows from $M\le b'\Sigma b$.

Write $R_n=dP/dQ|_{\F_n}$ for these kernels.
Under $P$, $Z_t=w_t'A_t/\sqrt{v_t}$ has conditional standard Gaussian law.
Choose a fixed unit whitened direction in the range of $\Sigma$ when $v_t=0$; its coefficient
$a_t$ is then zero. The likelihood-ratio increment is exactly
\eqref{eq:vanishing-lr} with this $a_t,Z_t$, while the linear factor is
$1-a_tZ_t+a_t^2Z_t^2$. The previous moment estimates provide bounded $a_t$,
$\sum_t\E|a_t|^3<\infty$ and $\sum_t\E a_t^4<\infty$.
Thus Appendix~\ref{app:vanishing}'s likelihood and summable-discrepancy
arguments apply without alteration, with $\sum_t v_t$ as the predictable
normalizer. They give $\log R_n/\log n\to\chi/8$, convergence of
$\log G_n-\log R_n$ and uniformly bounded expected absolute discrepancy.
The score martingale $J_n-C_n$ has conditional variance
$v_t+\gamma^2v_t^2/2$, at most a fixed multiple of $v_t$.
The same martingale normalization gives $J_n/C_n\to1$.

The fixed-e-process upper bound, the change-of-measure bound for
level-indexed stopping rules, and the first-passage inversion in
Appendix~\ref{app:vanishing} now apply with $\kappa=\chi/8$.
Since $C_n/\log n\to\chi/(2\gamma)$, their ratio gives the attained and
probability-lower-bound coefficient $4/\gamma$. All bounds use the same
adaptive null, not a null law selected anew at each confidence level.
This proves Theorem~\ref{thm:multivariate-waiting}.

\subsection{Geometric and constant certificates}
For Corollary~\ref{cor:multivariate-geometric-bet}, the public contracts imply
$\E(Y_t^2\mid\F_{t-1})\le\nu_t^2$ by Minkowski. Under $P$,
$\nu_t^2=s_t+O(s_t^{3/2}+s_t^2)$, with summable expected remainder.
For each fixed rule with a finite deterministic stake bound, the Catoni
expansion and martingale normalization in Appendix~\ref{app:vanishing} give
\[
 \log E_n^{\mathrm{geo},\lambda}
 =\sum_{t\le n}\{\gamma\lambda_tv_t/2-\lambda_t^2s_t/2\}
                       +o(\log n)\quad\text{almost surely}.
\]
The quadratic strong law also gives
$\sum_{t\le n}s_t/\log n\to\chi_K/\gamma^2$.

We need one additional weighted Gaussian law. For standard Gaussian partial
sums, put $X_m=S_m/\sqrt m$ and
$f(x)=(x'\mathsf A x)^2/(x'\mathsf B x)$, with value zero at a zero
denominator. Since $0\preceq\mathsf A\preceq\mathsf B$,
$0\le f(x)\le x'\mathsf A x$; thus $f$ is even and square integrable.
For $i\le j$, $X_i,X_j$ have cross-correlation $\sqrt{i/j}\,I_p$.
The centered Gaussian Hermite expansion of an even function has only terms
of total degree at least two. Orthogonality therefore gives
\[
 0\le\operatorname{Cov}\{f(X_i),f(X_j)\}
            \le(i/j)\operatorname{Var}\{f(Z)\}.
\]
Indeed each Hermite term of degree $k$ is multiplied by $(i/j)^{k/2}$.
It follows that $\sum_{m\le n}f(X_m)/m$ has mean
$\psi\log n+O(1)$ and variance $O(\log n)$. Positivity and the previous
subsequence argument prove its ratio to $\log n$ tends to $\psi$ almost
surely.

To pass to estimated targets, define
$F(w)=(w'\Sigma w)^2/(w'Kw)$ with the same zero convention. Where the
denominator is positive,
\[
 \nabla F(w)=
 \frac{4(w'\Sigma w)\Sigma w}{w'Kw}
 -\frac{2(w'\Sigma w)^2Kw}{(w'Kw)^2}.
\]
Because $w'\Sigma w\le w'Kw$, its norm is at most a fixed multiple of
$\norm w_2$. The function extends differentiably with zero gradient on
$\ker K$. Hence
$|F(w)-F(u)|\le C(\norm w_2+\norm u_2)\norm{w-u}_2$.
Inverse-curvature convergence, eventual inactivity of the constraint and
weighted Cesaro now yield
\[
 \frac1{\log n}\sum_{t\le n}\frac{v_t^2}{s_t}
                 \longrightarrow\frac{\psi}{\gamma^2},
\]
where zero-risk terms are zero. This argument works for singular covariance.

Completing the square pointwise gives
\[
 \gamma\lambda_tv_t/2-\lambda_t^2s_t/2
 =\frac{\gamma^2v_t^2}{8s_t}
     -\frac{s_t}{2}\left(\lambda_t-\frac{\gamma v_t}{2s_t}\right)^2
\]
when $s_t>0$. This proves the upper rate $\psi/8$. On the range of $K$,
the strong raw-moment law gives
\[
 \sup_{w'w>0,\ w'Kw>0}
 \frac{|w'(\widehat M_t-\Sigma)w|}{w'Kw}\to0
 \quad\text{almost surely}.
\]
The numerator vanishes on $\ker K$; use its inverse square root on its
range. Clipping at one is nonexpansive, so the observable stake differs from
$\gamma v_t/(2s_t)$ by $o(1)$ uniformly in directions. Since
$\sum s_t/\log n$ has a finite limit, its accumulated square-completion loss
is $o(\log n)$. It attains $\psi/8$. Constant stakes give the stated
quadratic rate. Cauchy--Schwarz yields
$\chi^2\le\psi\chi_K$, while $f(Z)\le Z'\mathsf A Z$ gives $\psi\le\chi$.
First-passage inversion and the utility limits prove the coefficient claims.

For deterministic evaluation of $\psi$, Gaussian integration and
$1/y=\int_0^\infty e^{-ty}\,dt$ give
\[
 \psi=\int_0^\infty
 \frac{2\operatorname{tr}\{(\mathsf A C_t)^2\}
       +\{\operatorname{tr}(\mathsf A C_t)\}^2}
      {\sqrt{\det(I_p+2t\mathsf B)}}\,dt,\qquad
 C_t=(I_p+2t\mathsf B)^{-1}.
\]
Tonelli applies to the nonnegative integrand; singular zero directions
contribute no numerator. In two isotropic dimensions, separating the
Gaussian radius and angle gives the displayed geometric-mean formula.

For Corollary~\ref{cor:multivariate-constant-certificate}, under $P$,
\[
 \E(Y_t^2\mid\F_{t-1})=v_t+3\gamma^2v_t^2/4
       \le V+3\gamma^2V^2/4,
 \qquad d_*:=\nu_*^2-V-3\gamma^2V^2/4
                 =\gamma\sqrt3\,V^{3/2}>0.
\]
The policy-independent compensator is
$\mathcal B_n=\sum_{t\le n}\gamma^2v_t^2/(8d_*)$. Its expected total is
at most the finite constant
\[
 B_*=
 \frac{\norm\Sigma_{\op}^2
       \{(\operatorname{tr}\Sigma)^2+2\operatorname{tr}(\Sigma^2)\}}
      {8\gamma^2r^4d_*}
       \sum_{m=n_0}^{\infty}m^{-2}.
\]
Completing the square in the conditional Catoni-factor expectation, then
Ville and Markov, gives $B_*/u+\alpha e^u$ exactly as in the proof of
Proposition~\ref{prop:vanishing-uniform}. Taking $u=\log(1/\alpha)/2$
proves convergence to zero. The bound may be numerically uninformative for
a very small positive ridge; no finite-level efficiency claim is deduced.
